\documentclass[10pt]{article}
% Hojas resumen: un A4 por delante y por detrás, dos columnas por cara.
% Uso: \documentclass[10pt]{article} % Hojas resumen: un A4 por delante y por detrás, dos columnas por cara.
% Uso: \documentclass[10pt]{article} % Hojas resumen: un A4 por delante y por detrás, dos columnas por cara.
% Uso: \documentclass[10pt]{article} \input{estilo-resumen.tex} \bloque{analisis|algebra|azar} \cabecera{n}{Título}{Lecciones}
\usepackage[utf8]{inputenc}
\usepackage[T1]{fontenc}
\usepackage[spanish,es-nodecimaldot]{babel}
\usepackage{lmodern}
\usepackage{amsmath,amssymb}
\usepackage{graphicx}
\usepackage[a4paper,left=10mm,right=10mm,top=15mm,bottom=13mm,headheight=11pt,headsep=3mm,footskip=6mm]{geometry}
\usepackage{multicol}
\setlength{\columnsep}{6mm}
\setlength{\columnseprule}{0.3pt}
\usepackage{xcolor}
\usepackage{enumitem}
\setlist{nosep,leftmargin=1.1em,itemsep=1pt}
\usepackage{array}
\usepackage{needspace}
\usepackage{pgfplots}
\usepgfplotslibrary{fillbetween}
\pgfplotsset{compat=1.18,
  mini/.style={width=4.3cm,height=3.3cm,scale only axis=false,axis lines=middle,xtick=\empty,ytick=\empty,samples=120,
    xlabel={\scriptsize$x$},ylabel={\scriptsize$y$},title style={font=\scriptsize,yshift=-4pt},clip=true,
    every axis plot/.append style={thick,color=acento}}}
\usepackage{ragged2e}
\usepackage[most]{tcolorbox}
\usepackage{fancyhdr,lastpage}
\usepackage[hidelinks,pdfencoding=auto]{hyperref}
\hypersetup{pdflang=es-ES,pdfsubject={Matemáticas Aplicadas a las Ciencias Sociales II}}

\definecolor{analisis}{HTML}{1D4ED8}
\definecolor{algebra}{HTML}{86198F}
\definecolor{azar}{HTML}{9A3412}
\definecolor{aviso}{HTML}{B91C1C}
\colorlet{acento}{analisis}
\newcommand{\bloque}[1]{\colorlet{acento}{#1}}

\setlength{\parindent}{0pt}
\setlength{\parskip}{2pt}
\renewcommand{\arraystretch}{1.2}
\AtBeginDocument{\RaggedRight\fontsize{9.8}{11.9}\selectfont\setlength{\abovedisplayskip}{3pt}\setlength{\belowdisplayskip}{3pt}\setlength{\abovedisplayshortskip}{1pt}\setlength{\belowdisplayshortskip}{1pt}}

\newcommand{\temanum}{}\newcommand{\tematitulo}{}
\pagestyle{fancy}
\fancyhf{}
\fancyhead[L]{\footnotesize\color{acento}\textbf{Hoja resumen}\quad Tema \temanum. \tematitulo}
\fancyhead[R]{\footnotesize Matemáticas Aplicadas a las Ciencias Sociales II --- 2.\textsuperscript{o} Bachillerato}
\fancyfoot[L]{\scriptsize IES Al-Ándalus --- Juan María Gámez Ortiz, con la ayuda de Claude y de Codex --- CC BY-NC 4.0}
\fancyfoot[R]{\scriptsize Cara \thepage\ de \pageref{LastPage}}
\renewcommand{\headrulewidth}{0.4pt}

% Cabecera de la primera cara: número, título y lecciones del tema.
\newcommand{\cabecera}[3]{%
  \renewcommand{\temanum}{#1}\renewcommand{\tematitulo}{#2}%
  \hypersetup{pdftitle={Hoja resumen. Tema #1. #2}}%
  \noindent\begin{tcolorbox}[enhanced,colback=acento!8,colframe=acento,boxrule=0.6pt,arc=2pt,left=4pt,right=4pt,top=2pt,bottom=2pt]
  \begin{minipage}[c]{1.1cm}\includegraphics[width=1cm]{logo-al-andalus.png}\end{minipage}%
  \begin{minipage}[c]{\dimexpr\linewidth-1.2cm}{\large\bfseries\color{acento}Tema #1. #2}\\[1pt]{\scriptsize #3}\end{minipage}
  \end{tcolorbox}\vspace{-2pt}}

% Apartado con filete del color del bloque.
\newcommand{\seccion}[1]{\par\needspace{5\baselineskip}\vspace{3pt}{\color{acento}\bfseries\normalsize #1}\par\vspace{-1pt}{\color{acento}\hrule height 0.5pt}\vspace{2pt}}
\newcommand{\sub}[1]{\par\needspace{3\baselineskip}\vspace{1pt}\textbf{\color{acento}#1}\enspace}

% Fórmula clave, aviso de errores y ejemplo.
\tcbset{before upper=\RaggedRight}
\newtcolorbox{formula}{enhanced,colback=acento!6,colframe=acento!70,boxrule=0.5pt,arc=1.5pt,left=3pt,right=3pt,top=0pt,bottom=0pt,before skip=2pt,after skip=2pt}
\newtcolorbox{ojo}{enhanced,colback=aviso!5,colframe=aviso!60,boxrule=0.5pt,arc=1.5pt,left=3pt,right=3pt,top=1pt,bottom=1pt,before skip=3pt,after skip=2pt,
  fonttitle=\bfseries\footnotesize,coltitle=aviso,colbacktitle=aviso!5,title={Ojo},attach title to upper={\enspace}}
\newtcolorbox{ejemplo}{enhanced,colback=white,colframe=acento!45,boxrule=0.4pt,arc=1.5pt,left=3pt,right=3pt,top=1pt,bottom=1pt,before skip=3pt,after skip=2pt,
  fonttitle=\bfseries\footnotesize,coltitle=acento,colbacktitle=white,title={Ejemplo},attach title to upper={\enspace}}
\newcommand{\dx}{\,dx}
 \bloque{analisis|algebra|azar} \cabecera{n}{Título}{Lecciones}
\usepackage[utf8]{inputenc}
\usepackage[T1]{fontenc}
\usepackage[spanish,es-nodecimaldot]{babel}
\usepackage{lmodern}
\usepackage{amsmath,amssymb}
\usepackage{graphicx}
\usepackage[a4paper,left=10mm,right=10mm,top=15mm,bottom=13mm,headheight=11pt,headsep=3mm,footskip=6mm]{geometry}
\usepackage{multicol}
\setlength{\columnsep}{6mm}
\setlength{\columnseprule}{0.3pt}
\usepackage{xcolor}
\usepackage{enumitem}
\setlist{nosep,leftmargin=1.1em,itemsep=1pt}
\usepackage{array}
\usepackage{needspace}
\usepackage{pgfplots}
\usepgfplotslibrary{fillbetween}
\pgfplotsset{compat=1.18,
  mini/.style={width=4.3cm,height=3.3cm,scale only axis=false,axis lines=middle,xtick=\empty,ytick=\empty,samples=120,
    xlabel={\scriptsize$x$},ylabel={\scriptsize$y$},title style={font=\scriptsize,yshift=-4pt},clip=true,
    every axis plot/.append style={thick,color=acento}}}
\usepackage{ragged2e}
\usepackage[most]{tcolorbox}
\usepackage{fancyhdr,lastpage}
\usepackage[hidelinks,pdfencoding=auto]{hyperref}
\hypersetup{pdflang=es-ES,pdfsubject={Matemáticas Aplicadas a las Ciencias Sociales II}}

\definecolor{analisis}{HTML}{1D4ED8}
\definecolor{algebra}{HTML}{86198F}
\definecolor{azar}{HTML}{9A3412}
\definecolor{aviso}{HTML}{B91C1C}
\colorlet{acento}{analisis}
\newcommand{\bloque}[1]{\colorlet{acento}{#1}}

\setlength{\parindent}{0pt}
\setlength{\parskip}{2pt}
\renewcommand{\arraystretch}{1.2}
\AtBeginDocument{\RaggedRight\fontsize{9.8}{11.9}\selectfont\setlength{\abovedisplayskip}{3pt}\setlength{\belowdisplayskip}{3pt}\setlength{\abovedisplayshortskip}{1pt}\setlength{\belowdisplayshortskip}{1pt}}

\newcommand{\temanum}{}\newcommand{\tematitulo}{}
\pagestyle{fancy}
\fancyhf{}
\fancyhead[L]{\footnotesize\color{acento}\textbf{Hoja resumen}\quad Tema \temanum. \tematitulo}
\fancyhead[R]{\footnotesize Matemáticas Aplicadas a las Ciencias Sociales II --- 2.\textsuperscript{o} Bachillerato}
\fancyfoot[L]{\scriptsize IES Al-Ándalus --- Juan María Gámez Ortiz, con la ayuda de Claude y de Codex --- CC BY-NC 4.0}
\fancyfoot[R]{\scriptsize Cara \thepage\ de \pageref{LastPage}}
\renewcommand{\headrulewidth}{0.4pt}

% Cabecera de la primera cara: número, título y lecciones del tema.
\newcommand{\cabecera}[3]{%
  \renewcommand{\temanum}{#1}\renewcommand{\tematitulo}{#2}%
  \hypersetup{pdftitle={Hoja resumen. Tema #1. #2}}%
  \noindent\begin{tcolorbox}[enhanced,colback=acento!8,colframe=acento,boxrule=0.6pt,arc=2pt,left=4pt,right=4pt,top=2pt,bottom=2pt]
  \begin{minipage}[c]{1.1cm}\includegraphics[width=1cm]{logo-al-andalus.png}\end{minipage}%
  \begin{minipage}[c]{\dimexpr\linewidth-1.2cm}{\large\bfseries\color{acento}Tema #1. #2}\\[1pt]{\scriptsize #3}\end{minipage}
  \end{tcolorbox}\vspace{-2pt}}

% Apartado con filete del color del bloque.
\newcommand{\seccion}[1]{\par\needspace{5\baselineskip}\vspace{3pt}{\color{acento}\bfseries\normalsize #1}\par\vspace{-1pt}{\color{acento}\hrule height 0.5pt}\vspace{2pt}}
\newcommand{\sub}[1]{\par\needspace{3\baselineskip}\vspace{1pt}\textbf{\color{acento}#1}\enspace}

% Fórmula clave, aviso de errores y ejemplo.
\tcbset{before upper=\RaggedRight}
\newtcolorbox{formula}{enhanced,colback=acento!6,colframe=acento!70,boxrule=0.5pt,arc=1.5pt,left=3pt,right=3pt,top=0pt,bottom=0pt,before skip=2pt,after skip=2pt}
\newtcolorbox{ojo}{enhanced,colback=aviso!5,colframe=aviso!60,boxrule=0.5pt,arc=1.5pt,left=3pt,right=3pt,top=1pt,bottom=1pt,before skip=3pt,after skip=2pt,
  fonttitle=\bfseries\footnotesize,coltitle=aviso,colbacktitle=aviso!5,title={Ojo},attach title to upper={\enspace}}
\newtcolorbox{ejemplo}{enhanced,colback=white,colframe=acento!45,boxrule=0.4pt,arc=1.5pt,left=3pt,right=3pt,top=1pt,bottom=1pt,before skip=3pt,after skip=2pt,
  fonttitle=\bfseries\footnotesize,coltitle=acento,colbacktitle=white,title={Ejemplo},attach title to upper={\enspace}}
\newcommand{\dx}{\,dx}
 \bloque{analisis|algebra|azar} \cabecera{n}{Título}{Lecciones}
\usepackage[utf8]{inputenc}
\usepackage[T1]{fontenc}
\usepackage[spanish,es-nodecimaldot]{babel}
\usepackage{lmodern}
\usepackage{amsmath,amssymb}
\usepackage{graphicx}
\usepackage[a4paper,left=10mm,right=10mm,top=15mm,bottom=13mm,headheight=11pt,headsep=3mm,footskip=6mm]{geometry}
\usepackage{multicol}
\setlength{\columnsep}{6mm}
\setlength{\columnseprule}{0.3pt}
\usepackage{xcolor}
\usepackage{enumitem}
\setlist{nosep,leftmargin=1.1em,itemsep=1pt}
\usepackage{array}
\usepackage{needspace}
\usepackage{pgfplots}
\usepgfplotslibrary{fillbetween}
\pgfplotsset{compat=1.18,
  mini/.style={width=4.3cm,height=3.3cm,scale only axis=false,axis lines=middle,xtick=\empty,ytick=\empty,samples=120,
    xlabel={\scriptsize$x$},ylabel={\scriptsize$y$},title style={font=\scriptsize,yshift=-4pt},clip=true,
    every axis plot/.append style={thick,color=acento}}}
\usepackage{ragged2e}
\usepackage[most]{tcolorbox}
\usepackage{fancyhdr,lastpage}
\usepackage[hidelinks,pdfencoding=auto]{hyperref}
\hypersetup{pdflang=es-ES,pdfsubject={Matemáticas Aplicadas a las Ciencias Sociales II}}

\definecolor{analisis}{HTML}{1D4ED8}
\definecolor{algebra}{HTML}{86198F}
\definecolor{azar}{HTML}{9A3412}
\definecolor{aviso}{HTML}{B91C1C}
\colorlet{acento}{analisis}
\newcommand{\bloque}[1]{\colorlet{acento}{#1}}

\setlength{\parindent}{0pt}
\setlength{\parskip}{2pt}
\renewcommand{\arraystretch}{1.2}
\AtBeginDocument{\RaggedRight\fontsize{9.8}{11.9}\selectfont\setlength{\abovedisplayskip}{3pt}\setlength{\belowdisplayskip}{3pt}\setlength{\abovedisplayshortskip}{1pt}\setlength{\belowdisplayshortskip}{1pt}}

\newcommand{\temanum}{}\newcommand{\tematitulo}{}
\pagestyle{fancy}
\fancyhf{}
\fancyhead[L]{\footnotesize\color{acento}\textbf{Hoja resumen}\quad Tema \temanum. \tematitulo}
\fancyhead[R]{\footnotesize Matemáticas Aplicadas a las Ciencias Sociales II --- 2.\textsuperscript{o} Bachillerato}
\fancyfoot[L]{\scriptsize IES Al-Ándalus --- Juan María Gámez Ortiz, con la ayuda de Claude y de Codex --- CC BY-NC 4.0}
\fancyfoot[R]{\scriptsize Cara \thepage\ de \pageref{LastPage}}
\renewcommand{\headrulewidth}{0.4pt}

% Cabecera de la primera cara: número, título y lecciones del tema.
\newcommand{\cabecera}[3]{%
  \renewcommand{\temanum}{#1}\renewcommand{\tematitulo}{#2}%
  \hypersetup{pdftitle={Hoja resumen. Tema #1. #2}}%
  \noindent\begin{tcolorbox}[enhanced,colback=acento!8,colframe=acento,boxrule=0.6pt,arc=2pt,left=4pt,right=4pt,top=2pt,bottom=2pt]
  \begin{minipage}[c]{1.1cm}\includegraphics[width=1cm]{logo-al-andalus.png}\end{minipage}%
  \begin{minipage}[c]{\dimexpr\linewidth-1.2cm}{\large\bfseries\color{acento}Tema #1. #2}\\[1pt]{\scriptsize #3}\end{minipage}
  \end{tcolorbox}\vspace{-2pt}}

% Apartado con filete del color del bloque.
\newcommand{\seccion}[1]{\par\needspace{5\baselineskip}\vspace{3pt}{\color{acento}\bfseries\normalsize #1}\par\vspace{-1pt}{\color{acento}\hrule height 0.5pt}\vspace{2pt}}
\newcommand{\sub}[1]{\par\needspace{3\baselineskip}\vspace{1pt}\textbf{\color{acento}#1}\enspace}

% Fórmula clave, aviso de errores y ejemplo.
\tcbset{before upper=\RaggedRight}
\newtcolorbox{formula}{enhanced,colback=acento!6,colframe=acento!70,boxrule=0.5pt,arc=1.5pt,left=3pt,right=3pt,top=0pt,bottom=0pt,before skip=2pt,after skip=2pt}
\newtcolorbox{ojo}{enhanced,colback=aviso!5,colframe=aviso!60,boxrule=0.5pt,arc=1.5pt,left=3pt,right=3pt,top=1pt,bottom=1pt,before skip=3pt,after skip=2pt,
  fonttitle=\bfseries\footnotesize,coltitle=aviso,colbacktitle=aviso!5,title={Ojo},attach title to upper={\enspace}}
\newtcolorbox{ejemplo}{enhanced,colback=white,colframe=acento!45,boxrule=0.4pt,arc=1.5pt,left=3pt,right=3pt,top=1pt,bottom=1pt,before skip=3pt,after skip=2pt,
  fonttitle=\bfseries\footnotesize,coltitle=acento,colbacktitle=white,title={Ejemplo},attach title to upper={\enspace}}
\newcommand{\dx}{\,dx}

\bloque{algebra}
\begin{document}
\cabecera{6}{Programación lineal}{Lecciones 6.1--6.5: inecuaciones lineales con dos incógnitas, sistemas de inecuaciones y región factible, formulación de un problema, método gráfico y problemas de programación lineal.}
\begin{multicols}{2}

\seccion{1. Inecuación lineal: un semiplano}
\begin{formula}
\[ax+by\le c,\qquad (a,b)\ne(0,0)\]
\end{formula}
La recta \(ax+by=c\) (\textbf{frontera}) divide el plano en dos semiplanos; la inecuación elige uno.
\begin{enumerate}
  \item \textbf{Dibuja la frontera}: continua con \(\le,\ge\) (se incluye); discontinua con \(<,>\) (se excluye).
  \item \textbf{Punto de prueba} que \textbf{no} esté en la frontera; normalmente \((0,0)\). Si la recta pasa por el origen, usa \((0,1)\) o \((1,0)\).
  \item Si el punto cumple, sombrea su lado; si no, el contrario.
\end{enumerate}
\begin{ejemplo}
\(x+2y\le10\): frontera por \((0,5)\) y \((10,0)\); \(0+0\le10\) es cierto \(\Rightarrow\) el semiplano del origen, con la recta incluida.
\(y\ge2x\): el origen está en la frontera; con \((0,1)\): \(1\ge0\), se sombrea por encima.
\end{ejemplo}
\sub{Casos.} \(x\le3\): frontera vertical. \(y>-2\): horizontal y discontinua. \(-2y\le4\iff y\ge-2\): al dividir entre un negativo se invierte la desigualdad.

\seccion{2. Sistema de inecuaciones: región factible}
La \textbf{región factible} es la intersección de todos los semiplanos: los puntos que cumplen \textbf{todas} las restricciones.
\sub{Vértices:} cruces de dos fronteras que cumplen todas las restricciones. Un cruce que incumple alguna no es vértice.
\begin{ejemplo}
\(x+y\le60,\ x\le40,\ x,y\ge0\): vértices \((0,0)\), \((40,0)\), \((40,20)\) y \((0,60)\). El cruce \((60,0)\) no es vértice: incumple \(x\le40\).
\end{ejemplo}
\begin{tabular}{@{}p{2.3cm}p{5.5cm}@{}}
\textbf{Región} & \textbf{Ejemplo}\\ \hline
Acotada & Polígono: cabe en un rectángulo\\
No acotada & \(x,y\ge0,\ x+y\ge2\)\\
Vacía & \(x+y\le2\) y \(x+y\ge5\)\\
Segmento & \(x+y\le4,\ x+y\ge4,\ x,y\ge0\)\\
Un punto & \(x\ge1,\ y\ge2,\ x+y\le3\): \((1,2)\)
\end{tabular}

Área cero no significa vacía: un segmento tiene infinitas soluciones.

\seccion{3. Plantear un problema}
\begin{formula}
\textbf{Variables} \(x,y\) (con unidades) \(\cdot\) \textbf{Función objetivo} \(Z=ax+by\) (maximizar o minimizar) \(\cdot\) \textbf{Restricciones} lineales \(\cdot\) \textbf{No negatividad} \(x\ge0,\ y\ge0\) si son cantidades.
\end{formula}
\begin{tabular}{@{}p{4.3cm}p{3.5cm}@{}}
\textbf{Enunciado} & \textbf{Restricción}\\ \hline
Como mínimo 100 en total & \(x+y\ge100\)\\
Como máximo 500\,euros en total & \(x+y\le500\)\\
El doble de A que de B & \(x=2y\)\\
Al menos el doble de A que de B & \(x\ge2y\)\\
A no supera a B en más de 20 & \(x-y\le20\)\\
B es al menos el 40\,\% del total & \(y\ge0{,}4(x+y)\)\\
\guillemotleft{}Dispone de\guillemotright{} 18 h & \(\ldots\le18\) (no \(=\))
\end{tabular}

\seccion{4. Método gráfico}
\begin{enumerate}
  \item \textbf{Plantear}: variables, tabla de datos, objetivo y restricciones.
  \item \textbf{Dibujar} la región factible.
  \item \textbf{Calcular los vértices} y comprobar que son factibles.
  \item \textbf{Evaluar} \(Z\) en cada vértice.
  \item \textbf{Decidir} e interpretar con unidades.
\end{enumerate}
\begin{formula}
\textbf{Teorema:} si la región es \textbf{no vacía, cerrada y acotada}, una función lineal alcanza su máximo y su mínimo en algún vértice.
\end{formula}
\sub{Rectas de nivel.} \(ax+by=k\) son rectas paralelas; al desplazarlas en el sentido en que crece \(Z\), la última posición que toca la región da el máximo.

\seccion{5. Ejemplo resuelto: sillas y mesas}
18 h de corte y 42 h de montaje. Silla: 2 h y 1 h, 30\,euros. Mesa: 2 h y 3 h, 50\,euros.
\[\text{Max }Z=30x+50y\ \text{ s.a. }\begin{cases}2x+2y\le18\\x+3y\le42\\x,y\ge0\end{cases}\]
\begin{center}
\begin{tikzpicture}
\begin{axis}[mini,width=5.6cm,height=4.4cm,xmin=-0.8,xmax=15,ymin=-0.8,ymax=15,xtick={9},ytick={9,14},tick label style={font=\scriptsize}]
\addplot[fill=acento!20,draw=acento,thick] coordinates {(0,0) (9,0) (0,9)} -- cycle;
\addplot[domain=0:14,dashed]{(42-x)/3};
\addplot[domain=0:15,densely dotted,black]{(450-30*x)/50};
\addplot[only marks,mark=*,mark size=1.6pt,color=acento] coordinates {(0,0) (9,0) (0,9)};
\node[font=\scriptsize,anchor=south west] at (axis cs:0,9) {\(C\)};
\node[font=\scriptsize,anchor=north west] at (axis cs:9,0) {\(B\)};
\end{axis}
\end{tikzpicture}
\end{center}
El montaje es \textbf{redundante}: \(x+3y\le3(x+y)\le27<42\). Vértices \(A(0,0)\), \(B(9,0)\), \(C(0,9)\):
\[Z(A)=0,\quad Z(B)=270,\quad Z(C)=450\]
\textbf{Máximo: 0 sillas y 9 mesas, 450\,euros} (la línea de puntos es la recta de nivel \(Z=450\)).

\seccion{6. Casos especiales}
\begin{itemize}
  \item \textbf{Óptimo múltiple}: dos vértices empatan; todo el lado que los une es óptimo. Max \(2x+2y\) con \(x+y\le6\): vale 12 en todo el segmento de \((6,0)\) a \((0,6)\).
  \item \textbf{Región vacía}: no hay solución.
  \item \textbf{No acotado}: con \(x,y\ge0,\ x+y\ge2\), \(Z=x+y\) no tiene máximo (crece sin límite), pero sí mínimo, 2.
  \item \textbf{Frontera excluida}: con \(0\le x<1\), \(0\le y\le1\), \(Z=x+y\) se acerca a 2 sin alcanzarlo: no hay máximo.
\end{itemize}

\seccion{7. Soluciones enteras}
Si las variables cuentan objetos, el óptimo debe ser entero. Si el óptimo del modelo continuo ya es entero, sirve. Si no, \textbf{no basta redondear}: hay que comprobar los puntos enteros cercanos en las restricciones y en \(Z\).

\seccion{8. Problemas tipo}
\begin{itemize}
  \item \textbf{Producción} (maximizar beneficio con recursos limitados): restricciones con \(\le\).
  \item \textbf{Dieta o mezcla} (minimizar coste con requisitos mínimos): restricciones con \(\ge\); la región suele ser no acotada, y el mínimo existe.
  \item \textbf{Inversión} (repartir un capital con límites y proporciones).
\end{itemize}
\seccion{9. Mínimo en una región no acotada: la dieta}
A cuesta 4 euros/kg y B, 3 euros/kg. Proteínas \(2x+4y\ge12\), vitaminas \(6x+2y\ge12\), \(x,y\ge0\). Minimizar \(Z=4x+3y\).
\begin{center}
\begin{tikzpicture}
\begin{axis}[mini,width=5.4cm,height=4.2cm,xmin=-0.5,xmax=7.5,ymin=-0.5,ymax=7.5]
\addplot[fill=acento!20,draw=none] coordinates {(0,7.5) (0,6) (1.2,2.4) (6,0) (7.5,0) (7.5,7.5)} -- cycle;
\addplot[domain=0:7.5]{(6-x)/2};
\addplot[domain=0:2.16]{6-3*x};
\addplot[only marks,mark=*,mark size=1.6pt,color=acento] coordinates {(0,6) (1.2,2.4) (6,0)};
\end{axis}
\end{tikzpicture}
\end{center}
Simplificando: \(x+2y\ge6\) y \(3x+y\ge6\). Vértices y valores:
\[\begin{array}{c|ccc}\text{Vértice}&(0,6)&(1{,}2;\ 2{,}4)&(6,0)\\\hline Z&18&\mathbf{12}&24\end{array}\]
La región no está acotada, pero el mínimo existe: \(Z=(x+2y)+(3x+y)\ge6+6=12\). \textbf{Mínimo: 1,2 kg de A y 2,4 kg de B, 12 euros.} No hay máximo.

\seccion{10. Cuatro vértices: la inversión}
50 mil euros como máximo; en A como mucho 30, en B al menos 15, y A al menos tanto como B. Interés: 10\,\% en A y 5\,\% en B.
\[\text{Max }Z=0{,}10x+0{,}05y\ \text{ s.a. }\begin{cases}x+y\le50,\ x\le30\\ y\ge15,\ x\ge y\end{cases}\]
\[\begin{array}{c|cccc}\text{Vértice}&(15,15)&(30,15)&(30,20)&(25,25)\\\hline Z&2{,}25&3{,}75&\mathbf{4}&3{,}75\end{array}\]
\textbf{Invertir 30 mil euros en A y 20 en B: 4 mil euros de interés anual.}

\seccion{11. Una restricción redundante: los muebles}
120 kg de madera y 100 h. Silla: 10 kg, 5 h, 20 euros. Mesa: 20 kg, 10 h, 35 euros.
\[\begin{cases}10x+20y\le120\iff x+2y\le12\\5x+10y\le100\iff x+2y\le20\end{cases}\]
La segunda es redundante: si \(x+2y\le12\), también \(x+2y\le20\). Vértices \((0,0)\), \((12,0)\), \((0,6)\) con \(Z=20x+35y\): 0, \textbf{240} y 210. \textbf{Máximo: 12 sillas y ninguna mesa, 240 euros.}

\seccion{12. Comprobar un punto}
Un punto es factible si cumple \textbf{todas} las restricciones; basta que falle una para descartarlo. En sillas y mesas:
\begin{itemize}
  \item \((3,4)\): corte \(2\cdot3+2\cdot4=14\le18\), montaje \(3+12=15\le42\). Factible, \(Z=290\) euros.
  \item \((6,4)\): corte \(20>18\). No factible, aunque su \(Z\) fuera mayor.
\end{itemize}

\seccion{13. Lista de control}
\begin{enumerate}
  \item \textquestiondown{}Variables definidas con unidades?
  \item \textquestiondown{}Cada frase del enunciado traducida a una restricción, con el sentido correcto?
  \item \textquestiondown{}No negatividad, si son cantidades?
  \item \textquestiondown{}Vértices comprobados en todas las restricciones?
  \item \textquestiondown{}Región acotada? Si no, \textquestiondown{}existe el óptimo pedido?
  \item \textquestiondown{}La respuesta dice cuánto de cada cosa y el valor de \(Z\) con unidades?
\end{enumerate}

\begin{ojo}
\begin{itemize}
  \item Usar el origen como punto de prueba cuando está en la frontera.
  \item Escribir igualdades donde el enunciado da máximos (\guillemotleft{}dispone de\guillemotright{}).
  \item Evaluar \(Z\) en cruces que no son factibles.
  \item Dar un solo punto cuando el óptimo es todo un lado.
  \item Concluir que una región no acotada no tiene óptimo.
\end{itemize}
\end{ojo}

\end{multicols}
\end{document}
